\documentclass[11pt]{article}
\usepackage{amsmath,amssymb,amsthm}
\newtheorem{theorem}{Theorem}
\newtheorem{proposition}[theorem]{Proposition}
\begin{document}
\title{Partial Fractions in the Boolean Quotient}
\author{}
\date{November 24, 2026}
\maketitle

\paragraph{Status.}
These are complete written algebraic identities and counterexamples,
not Lean-checked results. They concern exact cancellation using affine
numerators in the original admissible affine-ratio model. They do not
give an asymptotic head-complexity improvement. The counterexample fixes
the two denominators; it is not a lower bound against other denominators.

All equalities below are on the Boolean cube, or equivalently in the
quotient by $x_i^2-x_i$.

\section{A selector-bit product identity}

\begin{proposition}
Let $D_1,D_2$ be admissible affine denominators. Suppose
$D_2-D_1=c_0+c_1x_k$, where $c_0(c_0+c_1)\ne0$. Then the reciprocal
product is the sum of two admissible affine ratios:
\[
 \frac1{D_1D_2}=\frac{r(x_k)}{D_1}-\frac{r(x_k)}{D_2},\qquad
 r(x_k)=\frac{1-x_k}{c_0}+\frac{x_k}{c_0+c_1}.
\]
\end{proposition}
\begin{proof}
At either Boolean value of $x_k$, $r(x_k)(D_2-D_1)=1$.
Subtract the two fractions. The numerator $r$ is affine, and the
denominators have not changed.
\end{proof}

More generally, let $D_h=D+a_h+b_hx_k$ for $1\le h\le H$.
Suppose that for each $t\in\{0,1\}$ the numbers $a_h+b_ht$ are
pairwise distinct and every $D_h$ is admissible. Ordinary partial
fractions in $D$, applied separately at the two values of $x_k$, give
\[
 \prod_{h=1}^H\frac1{D_h}
 =\sum_{h=1}^H\frac{r_h(x_k)}{D_h},\qquad
 r_h(t)=\prod_{j\ne h}
            (a_j+b_jt-a_h-b_ht)^{-1}.
\]
Each $r_h$ has an affine interpolation between its two displayed
values. Thus a product of $H$ reciprocals in this family flattens
into $H$ heads exactly. If the offsets depend on several selector
bits, their partial-fraction residues are in general functions of
all of those bits, rather than affine numerators. The stated proof
therefore gives no constant-cost multibit extension.

There is a second small Boolean phenomenon. For $z_i,z_j\in\{-1,1\}$
and $a^2\ne b^2$,
\[
 (az_i+bz_j)^{-1}=\frac{az_i-bz_j}{a^2-b^2}.
\]
Consequently a difference of denominators of this centered two-bit
form also has an affine reciprocal. In $0,1$ coordinates, the
mixed finite difference of $1/(c+ax+by)$ is
\[
 \frac{ab(2c+a+b)}{c(c+a)(c+b)(c+a+b)}.
\]
Provided all four denominator values are nonzero, the reciprocal is
affine exactly when $ab(2c+a+b)=0$. These signed differences need
not themselves be admissible denominators; only the original
$D_h$ require admissibility.

\section{A genuinely multibit exact cancellation}

\begin{theorem}
On $\{0,1\}^3$, define
\[
\begin{split}
 D_1&=1+x+2y+3z,\qquad D_2=2+2x+5y+10z,\\
 N_1&=-\frac{33}{10}-\frac{x}{20}-\frac{37y}{20}
                                      +\frac{17z}{20},\\
 N_2&=\frac{38}{5}-\frac{2x}{5}+\frac{19y}{4}.
\end{split}
\]
Then
\[
             \frac1{D_1D_2}=\frac{N_1}{D_1}+\frac{N_2}{D_2}.
\]
Both heads are admissible on the entire cube.
\end{theorem}
\begin{proof}
Expansion before reduction gives
\[
 N_1D_2+N_2D_1
 =1+\frac12(x-x^2)-\frac14(y-y^2)-\frac{17}{2}(z-z^2).
\]
The three extra terms vanish on the cube. Both denominators have
positive constant terms and strictly positive slopes.
\end{proof}

This cancellation uses the Boolean relations rather than the
one-variable partial-fraction formula. The space of multilinear
polynomials of degree at most two in three bits has dimension seven,
while two affine numerators offer eight coefficients. Dimension alone
does not guarantee solvability for every denominator pair; the
displayed coefficients are the proof for this pair.

\section{A four-bit obstruction for fixed denominators}

\begin{theorem}
Let
\[
 D_1=1+\sum_{i=1}^4 i x_i,\qquad
 D_2=2+\sum_{i=1}^4(i^2+1)x_i.
\]
There do not exist affine $N_1,N_2$ with
$1/(D_1D_2)=N_1/D_1+N_2/D_2$ on $\{0,1\}^4$.
\end{theorem}
\begin{proof}
Both denominators are everywhere positive. The proposed identity
would therefore imply $N_1D_2+N_2D_1=1$ in the Boolean quotient.
On the space of degree-at-most-two multilinear polynomials, take
the coefficient order
\[
 (1,x_1,x_2,x_3,x_4,x_1x_2,x_1x_3,x_1x_4,
                          x_2x_3,x_2x_4,x_3x_4).
\]
Define a linear functional by the coefficient vector
\[
 \ell=(109,-92,-18,9,-2,90,36,-26,-36,18,0).
\]
Direct multiplication gives
\[
 \ell(D_1)=\ell(D_2)=0,\qquad
 \ell(x_iD_1)=\ell(x_iD_2)=0\quad(1\le i\le4),
\]
where products are reduced by $x_i^2=x_i$. Thus $\ell$ annihilates
every $N_1D_2+N_2D_1$ with affine numerators, but $\ell(1)=109$.
This is a contradiction.
\end{proof}

\paragraph{Exact remaining gap.}
The three-bit identity supplies a concrete cancellation mechanism,
but independent repetition gives only a constant improvement. A
useful global construction would have to couple many such
identities while keeping the resulting residues affine, or prove
a sign-only identity for a whole sum. None of the identities here
provides that coupling.
\end{document}
